% GENERATED FILE. Edit canonical lessons, then run npm run generate:books.
% canonical-source-hash: fnv1a64:bfc74c39ced96be4
% canonical-lessons: RU-C323-voyna, RU-C323-bitva, RU-C323-srazheniye, RU-C323-srazhatsya, RU-C323-vrag

\chapter{War}
\label{ch:ru-war}
\hlchaptermodality{323}

\begin{chapteropening}
\textbf{By the end of this chapter:} \emph{I can talk about a war, a battle, an engagement, fighting and an enemy.}

Complete the last lesson of chapter 323: I can talk about a war, a battle, an engagement, fighting and an enemy.
\end{chapteropening}

\section[\texorpdfstring{\ru{война} (voyná)}{война (voyná)}]{\ru{война} (voyná) — a war}

\label{lesson:RU-C323-voyna}



\begin{quote}
Before the new one: say the Russian for a shield, then the Russian for a throne.
\end{quote}

\subsection*{You'll want to know: \ru{война}}
\textbf{\ru{война}} (\emph{voyná}) — \textquotedblleft{}a war\textquotedblright{}.

\begin{quote}
\ru{Война} \ru{продолжалась} \ru{пять} \ru{лет}.
\end{quote}

\begin{quote}
\emph{Voyná prodolzhálas pyat let.}
\end{quote}

\begin{quote}
The war went on for five years.
\end{quote}

Feminine; plural \emph{\ru{войны}}. \emph{\ru{На} \ru{войне}} — in the war.

\begin{cousinweb}
\emph{\ru{война}} is inherited through Old East Slavic \emph{voina} from Proto-Slavic \emph{vojĭna} \textquotedblleft{}war\textquotedblright{}, built on \emph{vojĭ} \textquotedblleft{}a soldier\textquotedblright{} (EtymDB).
\end{cousinweb}

\begin{grammarlens}[title={What you've built}]
One more word for telling, explaining and describing at B1.
\end{grammarlens}

\subsection*{Your turn}
\begin{itemize}
\raggedright
  \item \emph{Say it:} \emph{voyná}
  \item \emph{Say it:} \emph{voyná}, once more
  \item \emph{Say it:} a sentence of your own with \emph{voyná}
  \item \emph{Recall:} say the Russian for a shield, then the Russian for a throne, then say \emph{voyná} again
\end{itemize}

\begin{tcolorbox}[breakable,colback=teal!4,colframe=teal!35!black,title={Before you move on}]
What does \ru{война} mean? (\textquotedblleft{}A war\textquotedblright{}.) Say it once more.
\end{tcolorbox}
\section[\texorpdfstring{\ru{битва} (bítva)}{битва (bítva)}]{\ru{битва} (bítva) — a battle}

\label{lesson:RU-C323-bitva}



\begin{quote}
Before the new one: say the Russian for a throne, then the Russian for a war.
\end{quote}

\subsection*{You'll want to know: \ru{битва}}
\textbf{\ru{битва}} (\emph{bítva}) — \textquotedblleft{}a battle\textquotedblright{}.

\begin{quote}
\ru{Битва} \ru{была} \ru{долгой} \ru{и} \ru{трудной}.
\end{quote}

\begin{quote}
\emph{Bítva bylá dólgoy i trúdnoy.}
\end{quote}

\begin{quote}
The battle was long and hard.
\end{quote}

Feminine.

\begin{cousinweb}
\emph{\ru{битва}} is inherited through Old East Slavic \emph{bitva} from Proto-Slavic \emph{bitva}, from PIE \emph{b\textsuperscript{h}ih\textsubscript{2}-tweh\textsubscript{2}} (EtymDB), the root of \emph{\ru{бить}} \textquotedblleft{}to beat\textquotedblright{}.
\end{cousinweb}

\begin{grammarlens}[title={What you've built}]
One more word for telling, explaining and describing at B1.
\end{grammarlens}

\subsection*{Your turn}
\begin{itemize}
\raggedright
  \item \emph{Say it:} \emph{bítva}
  \item \emph{Say it:} \emph{bítva}, once more
  \item \emph{Say it:} a sentence of your own with \emph{bítva}
  \item \emph{Recall:} say the Russian for a throne, then the Russian for a war, then say \emph{bítva} again
\end{itemize}

\begin{tcolorbox}[breakable,colback=teal!4,colframe=teal!35!black,title={Before you move on}]
What does \ru{битва} mean? (\textquotedblleft{}A battle\textquotedblright{}.) Say it once more.
\end{tcolorbox}
\section[\texorpdfstring{\ru{сражение} (srazhéniye)}{сражение (srazhéniye)}]{\ru{сражение} (srazhéniye) — a pitched battle, a battle}

\label{lesson:RU-C323-srazheniye}



\begin{quote}
Before the new one: say the Russian for a war, then the Russian for a battle.
\end{quote}

\subsection*{You'll want to know: \ru{сражение}}
\textbf{\ru{сражение}} (\emph{srazhéniye}) — \textquotedblleft{}a pitched battle, a battle\textquotedblright{}.

\begin{quote}
\ru{После} \ru{сражения} \ru{солдаты} \ru{отдыхали}.
\end{quote}

\begin{quote}
\emph{Pósle srazhéniya soldáty otdykháli.}
\end{quote}

\begin{quote}
After the battle the soldiers rested.
\end{quote}

Neuter. A little more formal than \emph{\ru{битва}}.

\begin{cousinweb}
\emph{\ru{сражение}} is built on \emph{\ru{сразиться}} \textquotedblleft{}to fight\textquotedblright{}, from \emph{\ru{с}-} and \emph{\ru{разить}} \textquotedblleft{}to strike\textquotedblright{} (EtymDB glosses \emph{\ru{разить}} \textquotedblleft{}to strike\textquotedblright{}); the \emph{\ru{з}} turns to \emph{\ru{ж}} in \emph{\ru{сражать}}, as in \emph{\ru{возить}}, \emph{\ru{вожу}}.
\end{cousinweb}

\begin{grammarlens}[title={What you've built}]
One more word for telling, explaining and describing at B1.
\end{grammarlens}

\subsection*{Your turn}
\begin{itemize}
\raggedright
  \item \emph{Say it:} \emph{srazhéniye}
  \item \emph{Say it:} \emph{srazhéniye}, once more
  \item \emph{Say it:} a sentence of your own with \emph{srazhéniye}
  \item \emph{Recall:} say the Russian for a war, then the Russian for a battle, then say \emph{srazhéniye} again
\end{itemize}

\begin{tcolorbox}[breakable,colback=teal!4,colframe=teal!35!black,title={Before you move on}]
What does \ru{сражение} mean? (\textquotedblleft{}A pitched battle, a battle\textquotedblright{}.) Say it once more.
\end{tcolorbox}
\section[\texorpdfstring{\ru{сражаться} (srazhátsya)}{сражаться (srazhátsya)}]{\ru{сражаться} (srazhátsya) — to fight}

\label{lesson:RU-C323-srazhatsya}



\begin{quote}
Before the new one: say the Russian for a battle, then the Russian for a pitched battle.
\end{quote}

\subsection*{You'll want to know: \ru{сражаться}}
\textbf{\ru{сражаться}} (\emph{srazhátsya}) — \textquotedblleft{}to fight\textquotedblright{}.

\begin{quote}
\ru{Солдаты} \ru{сражались} \ru{всю} \ru{ночь}.
\end{quote}

\begin{quote}
\emph{Soldáty srazhális vsyu noch.}
\end{quote}

\begin{quote}
The soldiers fought all night.
\end{quote}

Imperfective; perfective \emph{\ru{сразиться}} (\emph{srazítsya}). \emph{\ru{Сражаться} \ru{с}} + instrumental — against; \emph{\ru{за}} + accusative — for.

\begin{cousinweb}
\emph{\ru{сражаться}} is built on \emph{\ru{с}-} and \emph{\ru{разить}} \textquotedblleft{}to strike\textquotedblright{} (EtymDB), with the reflexive \emph{-\ru{ся}}: to strike at each other.
\end{cousinweb}

\begin{grammarlens}[title={What you've built}]
One more word for telling, explaining and describing at B1.
\end{grammarlens}

\subsection*{Your turn}
\begin{itemize}
\raggedright
  \item \emph{Say it:} \emph{srazhátsya}
  \item \emph{Say it:} \emph{srazhátsya}, once more
  \item \emph{Say it:} a sentence of your own with \emph{srazhátsya}
  \item \emph{Recall:} say the Russian for a battle, then the Russian for a pitched battle, then say \emph{srazhátsya} again
\end{itemize}

\begin{tcolorbox}[breakable,colback=teal!4,colframe=teal!35!black,title={Before you move on}]
What does \ru{сражаться} mean? (\textquotedblleft{}To fight\textquotedblright{}.) Say it once more.
\end{tcolorbox}
\section[\texorpdfstring{\ru{враг} (vrag)}{враг (vrag)}]{\ru{враг} (vrag) — an enemy}

\label{lesson:RU-C323-vrag}



\begin{quote}
Before the new one: say the Russian for a pitched battle, then the Russian for to fight.
\end{quote}

\subsection*{You'll want to know: \ru{враг}}
\textbf{\ru{враг}} (\emph{vrag}) — \textquotedblleft{}an enemy\textquotedblright{}.

\begin{quote}
\ru{Враг} \ru{подошёл} \ru{к} \ru{городу}.
\end{quote}

\begin{quote}
\emph{Vrag podoshyól k górodu.}
\end{quote}

\begin{quote}
The enemy approached the town.
\end{quote}

Genitive \emph{\ru{врага}}. The opposite of \emph{\ru{друг}}.

\begin{cousinweb}
\emph{\ru{враг}} is a Church Slavonic word, \emph{vragŭ} \textquotedblleft{}foe\textquotedblright{} (EtymDB), from Proto-Slavic \emph{vorgŭ} \textquotedblleft{}enemy\textquotedblright{}; the native Russian shape is the old \emph{\ru{ворог}}.
\end{cousinweb}

\begin{grammarlens}[title={What you've built}]
One more word for telling, explaining and describing at B1.
\end{grammarlens}

\subsection*{Your turn}
\begin{itemize}
\raggedright
  \item \emph{Say it:} \emph{vrag}
  \item \emph{Say it:} \emph{vrag}, once more
  \item \emph{Say it:} a sentence of your own with \emph{vrag}
  \item \emph{Recall:} say the Russian for a pitched battle, then the Russian for to fight, then say \emph{vrag} again
\end{itemize}

\begin{tcolorbox}[breakable,colback=teal!4,colframe=teal!35!black,title={Before you move on}]
What does \ru{враг} mean? (\textquotedblleft{}An enemy\textquotedblright{}.) Say it once more.
\end{tcolorbox}
